\section{Fichas e figuras de seleção da fase qualitativa}
\label{ap:selecao}

% O prefixo segue \thesection em vez da letra fixa: assim a numeração dos
% flutuantes acompanha a posição do apêndice se a ordem mudar, do mesmo
% modo que as referências cruzadas simbólicas já acompanham.
\setcounter{figure}{0}
\renewcommand{\thefigure}{\thesection.\arabic{figure}}
\setcounter{table}{0}
\renewcommand{\thetable}{\thesection.\arabic{table}}

Este apêndice reúne o material de apoio da seleção de entrevistados
descrita na Seção~\ref{sec:qualitativa}: as figuras de trajetória dos
hospitais escolhidos e as fichas do bloco de reserva. Foi
separado do corpo do texto por uma razão de proporção, são doze
figuras e uma lista de hospitais que ainda podem não ser
entrevistados, material de consulta e não de leitura corrida. As
decisões e os critérios continuam na Seção~\ref{sec:qualitativa}; aqui
está o detalhe que sustenta cada escolha.

\parte{Figuras de trajetória dos desviantes (blocos b e c)}

As sete figuras a seguir acompanham as fichas dos blocos b e c da
Seção~\ref{sec:qualitativa}.

\comoler{cada figura de seleção traz os quatro indicadores do hospital
(linha preta) contra a mediana da sua categoria (linha tracejada
colorida), de 2015 a 2025. Onde a linha preta corre abaixo da tracejada
em TMP, mortalidade e faturamento, ou acima em produção, o hospital supera a
referência da própria rede. Em Jundiaí, a produção corre bem acima da
mediana OSS e o TMP bem abaixo, a assinatura de um desviante positivo.}

\begin{figure}[H]
\centering
\includegraphics[width=.62\textwidth]{fig_sel_b_7573162.png}
\caption{Hospital Regional de Jundiaí (CNES 7573162), desviante positivo
da rede OSS.}
\label{fig:sel7573162}
\end{figure}

\comofeito{grade de quatro painéis, um para cada um dos quatro
indicadores da fase qualitativa: TMP (dias), mortalidade geral, faturamento
real por saída e produção (saídas hospitalares). Em cada painel, o eixo
horizontal é o ano (2015 a 2025) e o vertical, o indicador. A linha preta
sólida é o hospital e a linha tracejada colorida é a mediana da sua
categoria de gestão, para se ver onde o hospital supera ou fica abaixo da
própria rede.}

\begin{figure}[H]
\centering
\includegraphics[width=.62\textwidth]{fig_sel_b_6164366.png}
\caption{Hospital Estadual Américo Brasiliense (CNES 6164366), desviante
positivo da rede Direta.}
\label{fig:sel6164366}
\end{figure}

\comofeito{mesma construção da figura de seleção anterior: os quatro
indicadores (TMP, mortalidade, faturamento real e produção) do hospital em
preto contra a mediana tracejada da sua categoria, com o ano no eixo
horizontal.}

\begin{figure}[H]
\centering
\includegraphics[width=.62\textwidth]{fig_sel_b_5718368.png}
\caption{Hospital Municipal M'Boi Mirim (CNES 5718368), desviante
positivo da rede Público Municipal.}
\label{fig:sel5718368}
\end{figure}

\comofeito{mesma construção da figura de seleção anterior: os quatro
indicadores do hospital em preto contra a mediana tracejada da sua
categoria, ano no eixo horizontal.}

\begin{figure}[H]
\centering
\includegraphics[width=.62\textwidth]{fig_sel_b_2076896.png}
\caption{Hospital São Luiz Gonzaga (CNES 2076896), desviante positivo da
rede Filantrópico.}
\label{fig:sel2076896}
\end{figure}

\comofeito{mesma construção da figura de seleção anterior: os quatro
indicadores do hospital em preto contra a mediana tracejada da sua
categoria, ano no eixo horizontal.}

\comoler{mesma leitura das figuras do bloco anterior, agora ao
contrário: nestes hospitais a linha preta tende a correr acima da
mediana da categoria em TMP e faturamento, ou abaixo em produção, marcando a
posição no quartil inferior de eficiência.}

\begin{figure}[H]
\centering
\includegraphics[width=.62\textwidth]{fig_sel_c_2028840.png}
\caption{Instituto de Infectologia Emílio Ribas (CNES 2028840),
desviante negativo da rede Direta.}
\label{fig:sel2028840}
\end{figure}

\comofeito{mesma construção das figuras de seleção anteriores: os quatro
indicadores (TMP, mortalidade, faturamento real e produção) do hospital em
preto contra a mediana tracejada da sua categoria, ano no eixo
horizontal.}

\begin{figure}[H]
\centering
\includegraphics[width=.62\textwidth]{fig_sel_c_2698471.png}
\caption{Hospital Municipal Dr. Arthur Domingues Pinto (CNES 2698471),
desviante negativo da rede Público Municipal.}
\label{fig:sel2698471}
\end{figure}

\comofeito{mesma construção da figura de seleção anterior: os quatro
indicadores do hospital em preto contra a mediana tracejada da sua
categoria, ano no eixo horizontal.}

\begin{figure}[H]
\centering
\includegraphics[width=.62\textwidth]{fig_sel_c_2078465.png}
\caption{Real Sociedade Portuguesa de Beneficência (CNES 2078465),
desviante negativo da rede Filantrópico.}
\label{fig:sel2078465}
\end{figure}

\comofeito{mesma construção da figura de seleção anterior: os quatro
indicadores do hospital em preto contra a mediana tracejada da sua
categoria, ano no eixo horizontal.}

\parte{A reserva sugerida: entrevistas adicionais}

O bloco a seguir não integra o núcleo previsto; é uma
\textbf{sugestão} a acrescentar caso o tempo de campo permita. Os pares
casados aprofundam o contraste da comparação principal do documento.

\bloco{Bloco d: pares casados de contraste (sugeridos)}
\begin{center}\small
\begin{tabular}{llllr}
\toprule
CNES & Hospital & Município & Categoria & DEA \\
\midrule
2792141 & Hospital Regional De Cotia & Cotia & OSS & 0,893 \\
2079410 & Complexo Hospitalar Padre Bento De Gua & Guarulhos & Direta & 0,527 \\
2096412 & Santa Casa De Misericordia De Jacarei & Jacarei & Filantrópico & 0,598 \\
2065665 & Hospital Maternidade Interlagos & Sao Paulo & Direta & 0,690 \\
2079232 & Hospital Santa Barbara & Santa Barbara D'Oeste & Filantrópico & 0,523 \\
\bottomrule
\end{tabular}
\end{center}
\comofeito{mesmas cinco colunas da tabela do bloco a (CNES, hospital,
município, categoria e o escore mediano de DEA de 2015 a 2025); as linhas
alternam a âncora OSS e os seus pares casados de estrutura parecida.}

Duas âncoras OSS de alta eficiência, casadas cada uma a hospitais de
estrutura muito parecida mas eficiência bem menor. A ideia é entrevistar
os pares em espelho: mesma faixa Barcelona, mesmo porte, contagem de
leitos próxima, e ainda assim eficiência divergente. O que a gestão faz
de diferente?

\noindent\textbf{2792141, Hospital Regional de Cotia} (Cotia, Região Metropolitana de São Paulo; Médio Porte, faixa Barcelona 4, complexidade Barcelona 330,0). Âncora OSS de alta eficiência (DEA 0,893) para os pares casados.\par\medskip
\noindent\textbf{2079410, Complexo Hospitalar Padre Bento de Guarulhos} (Guarulhos, Região Metropolitana de São Paulo; Médio Porte, faixa Barcelona 4, complexidade Barcelona 327,0). Par casado da âncora OSS de Cotia (CNES 2792141): mesmo porte, faixa Barcelona 4, 145 leitos contra 145, mas eficiência divergente (DEA 0,527 contra 0,893). Estrutura parecida sob gestão Direta.\par\medskip
\noindent\textbf{2096412, Santa Casa de Misericórdia de Jacareí} (Jacareí, interior; Médio Porte, faixa Barcelona 4, complexidade Barcelona 312,0). Par casado da âncora OSS de Cotia (CNES 2792141): mesmo porte, faixa Barcelona 4, 143 leitos contra 145, mas eficiência divergente (DEA 0,598 contra 0,893). Estrutura parecida sob gestão Filantrópico.\par\medskip
\noindent\textbf{2065665, Hospital Maternidade Interlagos} (São Paulo, capital; Médio Porte, faixa Barcelona 3, complexidade Barcelona 160,0). Par casado da âncora OSS de Jundiaí (CNES 7573162): mesmo porte, faixa Barcelona 3, 89 leitos contra 136, mas eficiência divergente (DEA 0,690 contra 0,905). Estrutura parecida sob gestão Direta.\par\medskip
\noindent\textbf{2079232, Hospital Santa Bárbara} (Santa Bárbara d'Oeste, interior; Médio Porte, faixa Barcelona 3, complexidade Barcelona 247,0). Par casado da âncora OSS de Jundiaí (CNES 7573162): mesmo porte, faixa Barcelona 3, 137 leitos contra 136, mas eficiência divergente (DEA 0,523 contra 0,905). Estrutura parecida sob gestão Filantrópico.\par\medskip

\comoler{a leitura é a mesma das figuras anteriores. Nos pares casados,
o proveito está em comparar esta figura da âncora OSS com a do seu par
de estrutura parecida nas páginas seguintes: onde as trajetórias se
descolam, mora a diferença que a entrevista busca explicar.}

\begin{figure}[H]
\centering
\includegraphics[width=.62\textwidth]{fig_sel_d_2792141.png}
\caption{Hospital Regional de Cotia (CNES 2792141), âncora OSS de alta
eficiência dos pares casados.}
\label{fig:sel2792141}
\end{figure}

\comofeito{mesma construção das figuras de seleção anteriores: os quatro
indicadores do hospital em preto contra a mediana tracejada da sua
categoria, ano no eixo horizontal. Nos pares casados, o proveito está em
comparar esta âncora OSS com o seu par de estrutura parecida nas figuras
seguintes.}

\begin{figure}[H]
\centering
\includegraphics[width=.62\textwidth]{fig_sel_d_2079410.png}
\caption{Complexo Hospitalar Padre Bento de Guarulhos (CNES 2079410),
par Direta da âncora de Cotia.}
\label{fig:sel2079410}
\end{figure}

\comofeito{mesma construção da figura de seleção anterior: os quatro
indicadores do hospital em preto contra a mediana tracejada da sua
categoria, ano no eixo horizontal.}

\begin{figure}[H]
\centering
\includegraphics[width=.62\textwidth]{fig_sel_d_2096412.png}
\caption{Santa Casa de Misericórdia de Jacareí (CNES 2096412), par
Filantrópico da âncora de Cotia.}
\label{fig:sel2096412}
\end{figure}

\comofeito{mesma construção da figura de seleção anterior: os quatro
indicadores do hospital em preto contra a mediana tracejada da sua
categoria, ano no eixo horizontal.}

\begin{figure}[H]
\centering
\includegraphics[width=.62\textwidth]{fig_sel_d_2065665.png}
\caption{Hospital Maternidade Interlagos (CNES 2065665), par Direta da
âncora de Jundiaí.}
\label{fig:sel2065665}
\end{figure}

\comofeito{mesma construção da figura de seleção anterior: os quatro
indicadores do hospital em preto contra a mediana tracejada da sua
categoria, ano no eixo horizontal.}

\begin{figure}[H]
\centering
\includegraphics[width=.62\textwidth]{fig_sel_d_2079232.png}
\caption{Hospital Santa Bárbara (CNES 2079232), par Filantrópico da
âncora de Jundiaí.}
\label{fig:sel2079232}
\end{figure}

\comofeito{mesma construção da figura de seleção anterior: os quatro
indicadores do hospital em preto contra a mediana tracejada da sua
categoria, ano no eixo horizontal.}
